\section{組換え生物の拡散防止・第一種第二種使用}

\subsection{平成25年度 問11}
\label{q:h25-reg-apt-11}

\noindent\textbf{原典：} sources/H25/aptitude.pdf、PDF物理12頁。

\textbf{公式正答：} 3（single）。

\paragraph{問題の要旨} 組換え生物の第二種使用で、衣服・放流・教育目的・廃棄を判定する。

\Needspace{4\baselineskip}\paragraph{確認した条件・記述（独自要約）}\mbox{}\par

\noindent\textbf{方向：} 個々の記述の○/×に適合する組合せ。\par

\noindent\textbf{並び：} ア・イ・ウ・エ\par

\noindent\textbf{ア：} 実験衣を着用し退出時には脱ぐ。 判定○：衣服を介した付着・移出を防ぐ適切な措置である。文科省Q3-7は必要な付着/感染防止措置を説明しており、全形態で一律の実験衣義務と読み替えない。\par

\noindent\textbf{イ：} 組換えオタマジャクシ・メダカを川へ放流する。 判定×：拡散防止された第二種使用の取扱いではない。動物はP1A等の措置を取り、野外利用は第一種の承認規程という別の手続を検討する。\par

\noindent\textbf{ウ：} 教育目的なら拡散防止措置は不要。 判定×：教育にも同じ法の適用がある。実験の目的だけで環境への拡散防止を除外できない。\par

\noindent\textbf{エ：} 不活化後の廃棄物を廃棄物関係法令等に従い処分。 判定○：不活化で生物学的危険を減らしても廃棄物の化学的・物理的性質に応じた処理の義務は消えない。\par

\paragraph{問われる核心} 拡散防止の使用区分と生物学的処置・廃棄物処理の両立

\paragraph{必要知識} 拡散防止の使用区分と生物学的処置・廃棄物処理の両立。アは○。衣服を介した付着・移出を防ぐ適切な措置である。文科省Q3-7は必要な付着/感染防止措置を説明しており、全形態で一律の実験衣義務と読み替えない。 イは×。拡散防止された第二種使用の取扱いではない。動物はP1A等の措置を取り、野外利用は第一種の承認規程という別の手続を検討する。 ウは×。教育にも同じ法の適用がある。実験の目的だけで環境への拡散防止を除外できない。 エは○。不活化で生物学的危険を減らしても廃棄物の化学的・物理的性質に応じた処理の義務は消えない。 よって○・×・×・○の候補3を選ぶ。

\paragraph{考え方・正答を選ぶ理由（査読待ち初稿）} アは○。衣服を介した付着・移出を防ぐ適切な措置である。文科省Q3-7は必要な付着/感染防止措置を説明しており、全形態で一律の実験衣義務と読み替えない。 イは×。拡散防止された第二種使用の取扱いではない。動物はP1A等の措置を取り、野外利用は第一種の承認規程という別の手続を検討する。 ウは×。教育にも同じ法の適用がある。実験の目的だけで環境への拡散防止を除外できない。 エは○。不活化で生物学的危険を減らしても廃棄物の化学的・物理的性質に応じた処理の義務は消えない。 よって○・×・×・○の候補3を選ぶ。

\Needspace{4\baselineskip}\paragraph{選択肢の検討（査読待ち初稿）}\mbox{}\par

\noindent\textbf{1：整合しない} ×・×・×・×。個別判断と異なるのはア・エである。各記述の理由は確認条件欄に示す。\par

\noindent\textbf{2：整合しない} ○・×・○・×。個別判断と異なるのはウ・エである。各記述の理由は確認条件欄に示す。\par

\noindent\textbf{3：整合する} ○・×・×・○。個別判断の○・×・×・○と全項目で一致する。\par

\noindent\textbf{4：整合しない} ○・×・×・×。個別判断と異なるのはエである。各記述の理由は確認条件欄に示す。\par

\noindent\textbf{5：整合しない} ×・○・○・○。個別判断と異なるのはア・イ・ウである。各記述の理由は確認条件欄に示す。\par

\paragraph{再利用する考え方} 活動・物質・制度の対象と評価の境界を固定し、用語の入替、全部/一部、法的義務/目標、出題時/現在を区別する。拡散防止の使用区分と生物学的処置・廃棄物処理の両立を個別の技術活動に結び付ける。

\paragraph{共通理論} \hyperref[theory:environment:biosafety]{第\ref*{theory:environment:biosafety}節}

\par 参考：\cite{env_biosafety_qa,env_biosafety_school,env_biosafety_law}。

\par 難易度 2（校準初稿）、優先度未評価。

\section{排出削減・オフセットとニュートラルの境界}

\subsection{令和元年度再試験 問11}
\label{q:r01-ret-apt-11}

\noindent\textbf{原典：} sources/R01\_retest/aptitude.pdf、PDF物理12頁。

\textbf{公式正答：} 1（single）。

\paragraph{問題の要旨} 温室効果ガス、低炭素社会、オフセットとニュートラルの定義を判定する。

\Needspace{4\baselineskip}\paragraph{確認した条件・記述（独自要約）}\mbox{}\par

\noindent\textbf{方向：} 個々の記述の○/×に適合する組合せ。\par

\noindent\textbf{並び：} ア・イ・ウ・エ\par

\noindent\textbf{ア：} 温室効果ガスはCO2・CH4・N2Oの3種類だけ。 判定×：HFC/PFC/SF6等もある。当初の京都の6種と自然の水蒸気を含む放射作用という別の集合を区別する。\par

\noindent\textbf{イ：} 暮らしの豊かさと環境保全が両立する低炭素社会を目指す。 判定○：排出を減らす社会構想であり、快適性の放棄又は無制限な排出の肯定ではない。\par

\noindent\textbf{ウ：} オフセットを残余排出全部の埋め合わせ済み状態とする。 判定×：この全量・状態の定義はニュートラルに対応する。オフセットは全部又は一部を埋め合わせる取組。\par

\noindent\textbf{エ：} ニュートラルを残余排出全部又は一部の埋め合わせ取組とする。 判定×：これはオフセットの定義であり、一部だけの相殺では対象全体が中立とは言えない。\par

\paragraph{問われる核心} 排出量会計の対象と全量/一部の相殺

\paragraph{必要知識} 排出量会計の対象と全量/一部の相殺。アは×。HFC/PFC/SF6等もある。当初の京都の6種と自然の水蒸気を含む放射作用という別の集合を区別する。 イは○。排出を減らす社会構想であり、快適性の放棄又は無制限な排出の肯定ではない。 ウは×。この全量・状態の定義はニュートラルに対応する。オフセットは全部又は一部を埋め合わせる取組。 エは×。これはオフセットの定義であり、一部だけの相殺では対象全体が中立とは言えない。 よって×・○・×・×の候補1を選ぶ。

\paragraph{考え方・正答を選ぶ理由（査読待ち初稿）} アは×。HFC/PFC/SF6等もある。当初の京都の6種と自然の水蒸気を含む放射作用という別の集合を区別する。 イは○。排出を減らす社会構想であり、快適性の放棄又は無制限な排出の肯定ではない。 ウは×。この全量・状態の定義はニュートラルに対応する。オフセットは全部又は一部を埋め合わせる取組。 エは×。これはオフセットの定義であり、一部だけの相殺では対象全体が中立とは言えない。 よって×・○・×・×の候補1を選ぶ。

\Needspace{4\baselineskip}\paragraph{選択肢の検討（査読待ち初稿）}\mbox{}\par

\noindent\textbf{1：整合する} ×・○・×・×。個別判断の×・○・×・×と全項目で一致する。\par

\noindent\textbf{2：整合しない} ×・×・○・○。個別判断と異なるのはイ・ウ・エである。各記述の理由は確認条件欄に示す。\par

\noindent\textbf{3：整合しない} ×・○・○・○。個別判断と異なるのはウ・エである。各記述の理由は確認条件欄に示す。\par

\noindent\textbf{4：整合しない} ○・○・×・×。個別判断と異なるのはアである。各記述の理由は確認条件欄に示す。\par

\noindent\textbf{5：整合しない} ○・○・○・○。個別判断と異なるのはア・ウ・エである。各記述の理由は確認条件欄に示す。\par

\paragraph{再利用する考え方} 活動・物質・制度の対象と評価の境界を固定し、用語の入替、全部/一部、法的義務/目標、出題時/現在を区別する。排出量会計の対象と全量/一部の相殺を個別の技術活動に結び付ける。

\paragraph{共通理論} \hyperref[theory:environment:carbon]{第\ref*{theory:environment:carbon}節}

\par 参考：\cite{env_kyoto2005,env_offset_guideline,env_neutral2011}。

\paragraph{関連問題} \hyperref[q:h30-reg-apt-13]{H30-REG-APT-13}, \hyperref[q:r06-reg-apt-05]{R06-REG-APT-05}

\par 難易度 2（校準初稿）、優先度未評価。

\subsection{令和6年度 問5}
\label{q:r06-reg-apt-05}

\noindent\textbf{原典：} sources/R06/aptitude.pdf、PDF物理5頁。

\textbf{公式正答：} 5（single）。

\par\textbf{検討未確定：} 公式5・初回5を維持。イの境界超過で回復不能変化が引き起こされるという強い表現を、原研究のリスク増大と区別。Oとする出題意図と科学的限定をAstra査読対象にする。

\paragraph{問題の要旨} 相殺の2用語、地球の限界、トップランナー基準を比較する。

\Needspace{4\baselineskip}\paragraph{確認した条件・記述（独自要約）}\mbox{}\par

\noindent\textbf{方向：} 個々の記述の○/×に適合する組合せ。\par

\noindent\textbf{並び：} ア・イ・ウ・エ\par

\noindent\textbf{ア：} ニュートラルを自ら削減後の全部又は一部のクレジット埋め合わせ取組とする。 判定×：一部相殺まで含むこの定義はオフセットである。全量の相殺済み状態を区別。\par

\noindent\textbf{イ：} 地球の限界の安全活動範囲を越すと自然資源の回復不能な変化が引き起こされる。 判定○：安全範囲と境界の概念は適切だが、一次研究は急激・不可逆変化のリスク増大と述べる。超過すれば必ず直ちに回復不能となる決定論には疑義。\par

\noindent\textbf{ウ：} 省エネ基準を商品化されている最良機器の性能以上に設定する。 判定○：最良性能を基礎とする目標である。2024原本の非化石エネルギー転換を含む法名を保持。すべての機器へ同じ数値を定める方式ではない。\par

\noindent\textbf{エ：} オフセットを残余排出の全部を埋め合わせた状態とする。 判定×：全量・状態のニュートラル定義との入替である。\par

\paragraph{問われる核心} 排出会計の全量条件と地球の限界・機器効率

\paragraph{必要知識} 排出会計の全量条件と地球の限界・機器効率。アは×。一部相殺まで含むこの定義はオフセットである。全量の相殺済み状態を区別。 イは○。安全範囲と境界の概念は適切だが、一次研究は急激・不可逆変化のリスク増大と述べる。超過すれば必ず直ちに回復不能となる決定論には疑義。 ウは○。最良性能を基礎とする目標である。2024原本の非化石エネルギー転換を含む法名を保持。すべての機器へ同じ数値を定める方式ではない。 エは×。全量・状態のニュートラル定義との入替である。 よって×・○・○・×の候補5を選ぶ。

\paragraph{考え方・正答を選ぶ理由（査読待ち初稿）} アは×。一部相殺まで含むこの定義はオフセットである。全量の相殺済み状態を区別。 イは○。安全範囲と境界の概念は適切だが、一次研究は急激・不可逆変化のリスク増大と述べる。超過すれば必ず直ちに回復不能となる決定論には疑義。 ウは○。最良性能を基礎とする目標である。2024原本の非化石エネルギー転換を含む法名を保持。すべての機器へ同じ数値を定める方式ではない。 エは×。全量・状態のニュートラル定義との入替である。 よって×・○・○・×の候補5を選ぶ。

\Needspace{4\baselineskip}\paragraph{選択肢の検討（査読待ち初稿）}\mbox{}\par

\noindent\textbf{1：整合しない} ○・○・○・○。個別判断と異なるのはア・エである。各記述の理由は確認条件欄に示す。\par

\noindent\textbf{2：整合しない} ×・○・○・○。個別判断と異なるのはエである。各記述の理由は確認条件欄に示す。\par

\noindent\textbf{3：整合しない} ○・×・×・○。個別判断と異なるのはア・イ・ウ・エである。各記述の理由は確認条件欄に示す。\par

\noindent\textbf{4：整合しない} ○・○・×・×。個別判断と異なるのはア・ウである。各記述の理由は確認条件欄に示す。\par

\noindent\textbf{5：整合する} ×・○・○・×。個別判断の×・○・○・×と全項目で一致する。\par

\paragraph{再利用する考え方} 活動・物質・制度の対象と評価の境界を固定し、用語の入替、全部/一部、法的義務/目標、出題時/現在を区別する。排出会計の全量条件と地球の限界・機器効率を個別の技術活動に結び付ける。

\paragraph{共通理論} \hyperref[theory:environment:carbon]{第\ref*{theory:environment:carbon}節}, \hyperref[theory:environment:planetary]{第\ref*{theory:environment:planetary}節}, \hyperref[theory:environment:toprunner]{第\ref*{theory:environment:toprunner}節}

\par 参考：\cite{env_neutral2011,env_offset_guideline,env_planetary,env_toprunner2019}。

\paragraph{関連問題} \hyperref[q:r01-ret-apt-11]{R01-RET-APT-11}

\par 難易度 2（校準初稿）、優先度未評価。

\section{循環資源・再使用再生利用と処理優先}

\subsection{平成24年度 問5}
\label{q:h24-reg-apt-05}

\noindent\textbf{原典：} sources/H24/aptitude.pdf、PDF物理4頁。

\textbf{公式正答：} 3（single）。

\paragraph{問題の要旨} 循環資源の五段階の処理優先順位を比較する。

\Needspace{4\baselineskip}\paragraph{確認した条件・記述（独自要約）}\mbox{}\par

\noindent\textbf{方向：} 最も適切な順\par

\paragraph{問われる核心} 資源価値を維持する利用段階と処理の一般優先

\paragraph{必要知識} 資源価値を維持する利用段階と処理の一般優先。第5〜7条の一般原則は発生抑制、再使用、再生利用、熱回収、適正処分であり3。第7条は技術的経済的可能性と環境負荷の低減を考慮する例外を設け、全物品に機械的に同じ処理を強制するものではない。排出者責任と拡大生産者責任は主体を別にした責任である。

\paragraph{考え方・正答を選ぶ理由（査読待ち初稿）} 第5〜7条の一般原則は発生抑制、再使用、再生利用、熱回収、適正処分であり3。第7条は技術的経済的可能性と環境負荷の低減を考慮する例外を設け、全物品に機械的に同じ処理を強制するものではない。排出者責任と拡大生産者責任は主体を別にした責任である。

\Needspace{4\baselineskip}\paragraph{選択肢の検討（査読待ち初稿）}\mbox{}\par

\noindent\textbf{1：整合しない} 再使用→発生抑制→再生利用→熱回収→適正処分。発生を避ける段階が再使用より後になっている。\par

\noindent\textbf{2：整合しない} 発生抑制→再使用→再生利用→適正処分→熱回収。熱回収可能なものを先に適正処分とする順が逆。\par

\noindent\textbf{3：整合する} 発生抑制→再使用→再生利用→熱回収→適正処分。発生を避け、形を保ち、材料を回収し、熱利用し、残りを処分する一般順位に一致。\par

\noindent\textbf{4：整合しない} 再使用→再生利用→発生抑制→熱回収→適正処分。発生抑制を既に発生した資源の利用より後に置く。\par

\noindent\textbf{5：整合しない} 再生利用→再使用→発生抑制→適正処分→熱回収。発生抑制・再使用・再生利用の順と熱回収/処分の順の双方が逆。\par

\paragraph{再利用する考え方} 活動・物質・制度の対象と評価の境界を固定し、用語の入替、全部/一部、法的義務/目標、出題時/現在を区別する。資源価値を維持する利用段階と処理の一般優先を個別の技術活動に結び付ける。

\paragraph{共通理論} \hyperref[theory:environment:circular]{第\ref*{theory:environment:circular}節}, \hyperref[theory:environment:responsibility]{第\ref*{theory:environment:responsibility}節}

\par 参考：\cite{env_circular_law,env_circular_outline}。

\paragraph{関連問題} \hyperref[q:r04-reg-apt-10]{R04-REG-APT-10}

\par 難易度 2（校準初稿）、優先度未評価。

\subsection{平成28年度 問14}
\label{q:h28-reg-apt-14}

\noindent\textbf{原典：} sources/H28/aptitude.pdf、PDF物理14頁。

\textbf{公式正答：} 5（single）。

\par\textbf{検討未確定：} 公式5・初回5を保持。エの小売引取義務は過去販売/同種買替条件が省略され、無条件義務と読むと不正確。2016年のテレビ対象と現在の有機EL追加を分ける。

\paragraph{問題の要旨} 基本法・廃棄物処理・管理票・家電の主体別責任を比較する。

\Needspace{4\baselineskip}\paragraph{確認した条件・記述（独自要約）}\mbox{}\par

\noindent\textbf{方向：} 個々の記述の○/×に適合する組合せ。\par

\noindent\textbf{並び：} ア・イ・ウ・エ\par

\noindent\textbf{ア：} 基本法と容器包装・家電・食品・建設資材等の個別法を整備。 判定○：基本原則と、製品ごとの回収処理制度という階層が異なる。\par

\noindent\textbf{イ：} 産業廃棄物は排出事業者の責任で適正に処理。 判定○：委託しても排出者の責任を無条件に移転できない。一般廃棄物/産業廃棄物は法の分類で、危険性の大小だけによらない。\par

\noindent\textbf{ウ：} 第三者委託時も最終処分まで必要な措置に努め管理票を使う。 判定○：管理票は運搬・中間処理・最終処分を確認する連絡記録である。費用を払っただけで管理が完了するわけではない。\par

\noindent\textbf{エ：} 家電4群で排出者が費用を負担し小売が引取り指定場所へ渡す。 判定○：制度概要は整合するが、小売の引取義務は過去販売したもの又は同種買替などの条件を持つ。原典は引取りを求められた対象機器とだけ記し条件を省略する。\par

\paragraph{問われる核心} 排出者責任と製品回収における主体・条件

\paragraph{必要知識} 排出者責任と製品回収における主体・条件。アは○。基本原則と、製品ごとの回収処理制度という階層が異なる。 イは○。委託しても排出者の責任を無条件に移転できない。一般廃棄物/産業廃棄物は法の分類で、危険性の大小だけによらない。 ウは○。管理票は運搬・中間処理・最終処分を確認する連絡記録である。費用を払っただけで管理が完了するわけではない。 エは○。制度概要は整合するが、小売の引取義務は過去販売したもの又は同種買替などの条件を持つ。原典は引取りを求められた対象機器とだけ記し条件を省略する。 よって○・○・○・○の候補5を選ぶ。

\paragraph{考え方・正答を選ぶ理由（査読待ち初稿）} アは○。基本原則と、製品ごとの回収処理制度という階層が異なる。 イは○。委託しても排出者の責任を無条件に移転できない。一般廃棄物/産業廃棄物は法の分類で、危険性の大小だけによらない。 ウは○。管理票は運搬・中間処理・最終処分を確認する連絡記録である。費用を払っただけで管理が完了するわけではない。 エは○。制度概要は整合するが、小売の引取義務は過去販売したもの又は同種買替などの条件を持つ。原典は引取りを求められた対象機器とだけ記し条件を省略する。 よって○・○・○・○の候補5を選ぶ。

\Needspace{4\baselineskip}\paragraph{選択肢の検討（査読待ち初稿）}\mbox{}\par

\noindent\textbf{1：整合しない} ○・○・○・×。個別判断と異なるのはエである。各記述の理由は確認条件欄に示す。\par

\noindent\textbf{2：整合しない} ○・○・×・○。個別判断と異なるのはウである。各記述の理由は確認条件欄に示す。\par

\noindent\textbf{3：整合しない} ○・×・○・○。個別判断と異なるのはイである。各記述の理由は確認条件欄に示す。\par

\noindent\textbf{4：整合しない} ×・○・○・○。個別判断と異なるのはアである。各記述の理由は確認条件欄に示す。\par

\noindent\textbf{5：整合する} ○・○・○・○。個別判断の○・○・○・○と全項目で一致する。\par

\paragraph{再利用する考え方} 活動・物質・制度の対象と評価の境界を固定し、用語の入替、全部/一部、法的義務/目標、出題時/現在を区別する。排出者責任と製品回収における主体・条件を個別の技術活動に結び付ける。

\paragraph{共通理論} \hyperref[theory:environment:responsibility]{第\ref*{theory:environment:responsibility}節}, \hyperref[theory:environment:appliance]{第\ref*{theory:environment:appliance}節}

\par 参考：\cite{env_circular_law,env_waste2003,env_appliance,env_appliance2024}。

\par 難易度 2（校準初稿）、優先度未評価。

\subsection{令和4年度 問10}
\label{q:r04-reg-apt-10}

\noindent\textbf{原典：} sources/R04/aptitude.pdf、PDF物理9頁。

\textbf{公式正答：} 4（single）。

\paragraph{問題の要旨} 循環型社会、循環的利用、再生利用、処理順位の定義を照合する。

\Needspace{4\baselineskip}\paragraph{確認した条件・記述（独自要約）}\mbox{}\par

\noindent\textbf{方向：} 個々の記述の○/×に適合する組合せ。\par

\noindent\textbf{並び：} ア・イ・ウ・エ\par

\noindent\textbf{ア：} 発生抑制・循環利用・適正処分で天然資源消費と環境負荷を低減。 判定○：第2条の循環型社会の目的構造に一致。\par

\noindent\textbf{イ：} 循環的利用は再使用・再生利用・熱回収。 判定○：利用の3方式と適正処分は別である。\par

\noindent\textbf{ウ：} 再生利用はそのまま製品又は部品として使用。 判定×：これは再使用の定義である。再生利用は原材料として利用する。\par

\noindent\textbf{エ：} 優先は抑制・再生利用・再使用・熱回収・処分。 判定×：一般順位の2位再使用と3位再生利用が入れ替わっている。\par

\paragraph{問われる核心} 製品・原材料・熱という資源利用形態と優先順位

\paragraph{必要知識} 製品・原材料・熱という資源利用形態と優先順位。アは○。第2条の循環型社会の目的構造に一致。 イは○。利用の3方式と適正処分は別である。 ウは×。これは再使用の定義である。再生利用は原材料として利用する。 エは×。一般順位の2位再使用と3位再生利用が入れ替わっている。 よって○・○・×・×の候補4を選ぶ。

\paragraph{考え方・正答を選ぶ理由（査読待ち初稿）} アは○。第2条の循環型社会の目的構造に一致。 イは○。利用の3方式と適正処分は別である。 ウは×。これは再使用の定義である。再生利用は原材料として利用する。 エは×。一般順位の2位再使用と3位再生利用が入れ替わっている。 よって○・○・×・×の候補4を選ぶ。

\Needspace{4\baselineskip}\paragraph{選択肢の検討（査読待ち初稿）}\mbox{}\par

\noindent\textbf{1：整合しない} ○・○・○・○。個別判断と異なるのはウ・エである。各記述の理由は確認条件欄に示す。\par

\noindent\textbf{2：整合しない} ×・○・×・○。個別判断と異なるのはア・エである。各記述の理由は確認条件欄に示す。\par

\noindent\textbf{3：整合しない} ○・×・○・×。個別判断と異なるのはイ・ウである。各記述の理由は確認条件欄に示す。\par

\noindent\textbf{4：整合する} ○・○・×・×。個別判断の○・○・×・×と全項目で一致する。\par

\noindent\textbf{5：整合しない} ○・×・○・○。個別判断と異なるのはイ・ウ・エである。各記述の理由は確認条件欄に示す。\par

\paragraph{再利用する考え方} 活動・物質・制度の対象と評価の境界を固定し、用語の入替、全部/一部、法的義務/目標、出題時/現在を区別する。製品・原材料・熱という資源利用形態と優先順位を個別の技術活動に結び付ける。

\paragraph{共通理論} \hyperref[theory:environment:circular]{第\ref*{theory:environment:circular}節}

\par 参考：\cite{env_circular_law}。

\paragraph{関連問題} \hyperref[q:h24-reg-apt-05]{H24-REG-APT-05}, \hyperref[q:r07-reg-apt-10]{R07-REG-APT-10}

\par 難易度 2（校準初稿）、優先度未評価。

\subsection{令和7年度 問10}
\label{q:r07-reg-apt-10}

\noindent\textbf{原典：} sources/R07/aptitude.pdf、PDF物理10頁。

\textbf{公式正答：} 3（single）。

\paragraph{問題の要旨} 循環資源、再使用、再生利用、熱回収、優先順位の5記述の適切数を求める。

\Needspace{4\baselineskip}\paragraph{確認した条件・記述（独自要約）}\mbox{}\par

\noindent\textbf{方向：} 適切な記述の数。\par

\noindent\textbf{ア：} 有価無価を問わない廃棄物等のうち有用なものを循環資源とする。○：法の対象を有価品だけに限定しない。\par

\noindent\textbf{イ：} 製品/部品として使うことを再生利用とする。×：再使用の定義。修理して利用することも再使用。\par

\noindent\textbf{ウ：} 原材料として利用することを再使用とする。×：再生利用の定義。\par

\noindent\textbf{エ：} 燃焼可能な資源を熱を得ることへ使う熱回収。○：材料回収とは利用形態が異なる。\par

\noindent\textbf{オ：} 抑制→再生利用→再使用→熱回収→処分。×：一般順位の再使用と再生利用が逆。\par

\paragraph{問われる核心} 循環資源の利用形態と法定一般順位

\paragraph{必要知識} 循環資源の利用形態と法定一般順位。ア○・イ×・ウ×・エ○・オ×で適切2件、候補3。再使用は製品や部品として形・機能を保つ利用、再生利用は原材料としての利用である。熱回収を加えた3利用と、発生抑制・適正処分を含む一般5段階を区別する。

\paragraph{考え方・正答を選ぶ理由（査読待ち初稿）} ア○・イ×・ウ×・エ○・オ×で適切2件、候補3。再使用は製品や部品として形・機能を保つ利用、再生利用は原材料としての利用である。熱回収を加えた3利用と、発生抑制・適正処分を含む一般5段階を区別する。

\Needspace{4\baselineskip}\paragraph{選択肢の検討（査読待ち初稿）}\mbox{}\par

\noindent\textbf{1：整合しない} 4。3誤りがあるため4件ではない。\par

\noindent\textbf{2：整合しない} 3。イウオを誤りとすると適切2件であり3件ではない。\par

\noindent\textbf{3：整合する} 2。アとエの2件で一致。\par

\noindent\textbf{4：整合しない} 1。正しいアエのどちらかを取り落とす。\par

\noindent\textbf{5：整合しない} 0。アエは法の定義と一致し0件ではない。\par

\paragraph{再利用する考え方} 活動・物質・制度の対象と評価の境界を固定し、用語の入替、全部/一部、法的義務/目標、出題時/現在を区別する。循環資源の利用形態と法定一般順位を個別の技術活動に結び付ける。

\paragraph{共通理論} \hyperref[theory:environment:circular]{第\ref*{theory:environment:circular}節}

\par 参考：\cite{env_circular_law}。

\paragraph{関連問題} \hyperref[q:r04-reg-apt-10]{R04-REG-APT-10}

\par 難易度 2（校準初稿）、優先度未評価。

\section{最適化・意思決定}

\subsection{平成24年度 問8}
\label{q:h24-reg-apt-08}

\noindent\textbf{原典：} sources/H24/aptitude.pdf、PDF物理7頁。

\textbf{公式正答：} 3（single）。

\par\textit{原典本文の校正と独自解説は未作成。分類のみの登録項目。}

\subsection{平成25年度 問15}
\label{q:h25-reg-apt-15}

\noindent\textbf{原典：} sources/H25/aptitude.pdf、PDF物理17頁。

\textbf{公式正答：} 3（single）。

\par\textit{原典本文の校正と独自解説は未作成。分類のみの登録項目。}

\subsection{平成28年度 問10}
\label{q:h28-reg-apt-10}

\noindent\textbf{原典：} sources/H28/aptitude.pdf、PDF物理10頁。

\textbf{公式正答：} 5（single）。

\par\textit{原典本文の校正と独自解説は未作成。分類のみの登録項目。}

\section{品質管理・統計}

\subsection{令和7年度 問2}
\label{q:r07-reg-apt-02}

\noindent\textbf{原典：} sources/R07/aptitude.pdf、PDF物理2頁。

\textbf{公式正答：} 1（single）。

\par\textit{原典本文の校正と独自解説は未作成。分類のみの登録項目。}

\section{環境法・循環・持続可能性}

\subsection{平成23年度 問14}
\label{q:h23-reg-apt-14}

\noindent\textbf{原典：} sources/H23/aptitude.pdf、PDF物理15頁。

\textbf{公式正答：} 5（single）。

\par\textbf{検討未確定：} 肢1の2011年までの取組企業数増加とISO普及・費用上昇との因果は一次時系列未確認。定義の検証と分ける。

\paragraph{問題の要旨} 最終処分ゼロ、生物濃縮、多様性、持続可能性、排出取引の5概念を区別する。

\Needspace{4\baselineskip}\paragraph{確認した条件・記述（独自要約）}\mbox{}\par

\noindent\textbf{方向：} 不適切\par

\paragraph{問われる核心} 資源循環・生態系・世代間公平と排出の埋め合わせの境界

\paragraph{必要知識} 資源循環・生態系・世代間公平と排出の埋め合わせの境界。削減証明や排出枠は埋め合わせに使える手段であるが、手段・制度の名称とオフセットという活動は同一でないため5が不適切。1の狭義の定義と企業数の時系列主張を分離し、後者を未確認として残す。

\paragraph{考え方・正答を選ぶ理由（査読待ち初稿）} 削減証明や排出枠は埋め合わせに使える手段であるが、手段・制度の名称とオフセットという活動は同一でないため5が不適切。1の狭義の定義と企業数の時系列主張を分離し、後者を未確認として残す。

\Needspace{4\baselineskip}\paragraph{選択肢の検討（査読待ち初稿）}\mbox{}\par

\noindent\textbf{1：定義適切・時系列要確認} 狭義のゼロエミッションは生産廃棄物の最終処分ゼロで、国内の取組企業が増加。。最終処分量を対象にする狭い定義は企業事例と整合。ISO14000は規格群である。2011時点の取組企業数増加とその因果の一次統計は未確認。\par

\noindent\textbf{2：適切（蓄積性物質）} 蓄積性物質が食物連鎖上位の生物で高濃度となる。。難分解・蓄積性の物質では捕食による取込と排出の差で上位濃度が上がり得る。すべての物質で必ず上昇するという主張ではない。\par

\noindent\textbf{3：適切} 多様性は種・遺伝子・生態系の3階層と時間変化を含む。。同じ種の中の遺伝的差、種の間、相互作用する生態系の差を別階層とする。静止した種数だけに限定しない。\par

\noindent\textbf{4：適切} 将来のニーズを満たす能力を損なわず現在のニーズを満たす開発。。WCEDの定義は現在の開発を認めつつ将来能力を保持する世代間条件である。\par

\noindent\textbf{5：不適切} オフセットは先進国間で取引できる排出削減証明や排出枠の取引制度。。これは京都の排出量取引・クレジットの説明である。オフセットは自らの削減後の残余を他所の削減吸収で埋め合わせる取組で、政府間取引に限定されない。\par

\paragraph{再利用する考え方} 活動・物質・制度の対象と評価の境界を固定し、用語の入替、全部/一部、法的義務/目標、出題時/現在を区別する。資源循環・生態系・世代間公平と排出の埋め合わせの境界を個別の技術活動に結び付ける。

\paragraph{共通理論} \hyperref[theory:environment:boundary]{第\ref*{theory:environment:boundary}節}, \hyperref[theory:environment:zero]{第\ref*{theory:environment:zero}節}, \hyperref[theory:environment:accumulation]{第\ref*{theory:environment:accumulation}節}, \hyperref[theory:environment:biodiversity]{第\ref*{theory:environment:biodiversity}節}, \hyperref[theory:environment:development]{第\ref*{theory:environment:development}節}, \hyperref[theory:environment:carbon]{第\ref*{theory:environment:carbon}節}

\par 参考：\cite{env_zero2008,env_foodchain,env_biodiversity,env_development,env_offset2008}。

\par 難易度 2（校準初稿）、優先度未評価。

\subsection{平成25年度 問10}
\label{q:h25-reg-apt-10}

\noindent\textbf{原典：} sources/H25/aptitude.pdf、PDF物理11頁。

\textbf{公式正答：} 3（single）。

\par\textit{原典本文の校正と独自解説は未作成。分類のみの登録項目。}

\subsection{平成28年度 問12}
\label{q:h28-reg-apt-12}

\noindent\textbf{原典：} sources/H28/aptitude.pdf、PDF物理12頁。

\textbf{公式正答：} 5（single）。

\par\textit{原典本文の校正と独自解説は未作成。分類のみの登録項目。}

\subsection{平成30年度 問13}
\label{q:h30-reg-apt-13}

\noindent\textbf{原典：} sources/H30/aptitude.pdf、PDF物理13頁。

\textbf{公式正答：} 5（single）。

\par\textbf{検討未確定：} 公式5・初回5を保持。エ後半の生活に不要な物質の濃縮を一律に異常・環境問題とする表現は、毒性・曝露・自然濃縮を省略し一般化に疑義。

\paragraph{問題の要旨} 排出の埋め合わせ、持続可能性、廃棄物ゼロ、生物濃縮を判定する。

\Needspace{4\baselineskip}\paragraph{確認した条件・記述（独自要約）}\mbox{}\par

\noindent\textbf{方向：} 個々の記述の○/×に適合する組合せ。\par

\noindent\textbf{並び：} ア・イ・ウ・エ\par

\noindent\textbf{ア：} オフセットは削減後の残余を削減活動への投資等で埋め合わせる。 判定○：自らの削減と他所の削減吸収という二段の区別に一致。\par

\noindent\textbf{イ：} 1987年WCED報告で将来と現在のニーズを共に考える。 判定○：将来世代がニーズを満たす能力を損なわないという世代間条件を読む。後代の需要を完全予測する要求ではない。\par

\noindent\textbf{ウ：} 産業間で廃棄物・副産物を資源化し社会の排出ゼロを目指す。 判定○：社会システムの構想であり、現実の全施設で既に完全な物質・エネルギー排出ゼロを達成したとの主張ではない。\par

\noindent\textbf{エ：} 外界より高濃度の蓄積、特に不要元素物質の濃縮は異常・環境問題。 判定○：前半は濃縮の説明。後半を一律に読むと強過ぎる。健康影響には物質の毒性・濃度・曝露期間が必要で、濃縮や不要という性質だけで被害は決まらない。\par

\paragraph{問われる核心} 資源と排出の境界、濃縮と有害性の区別

\paragraph{必要知識} 資源と排出の境界、濃縮と有害性の区別。アは○。自らの削減と他所の削減吸収という二段の区別に一致。 イは○。将来世代がニーズを満たす能力を損なわないという世代間条件を読む。後代の需要を完全予測する要求ではない。 ウは○。社会システムの構想であり、現実の全施設で既に完全な物質・エネルギー排出ゼロを達成したとの主張ではない。 エは○。前半は濃縮の説明。後半を一律に読むと強過ぎる。健康影響には物質の毒性・濃度・曝露期間が必要で、濃縮や不要という性質だけで被害は決まらない。 よって○・○・○・○の候補5を選ぶ。

\paragraph{考え方・正答を選ぶ理由（査読待ち初稿）} アは○。自らの削減と他所の削減吸収という二段の区別に一致。 イは○。将来世代がニーズを満たす能力を損なわないという世代間条件を読む。後代の需要を完全予測する要求ではない。 ウは○。社会システムの構想であり、現実の全施設で既に完全な物質・エネルギー排出ゼロを達成したとの主張ではない。 エは○。前半は濃縮の説明。後半を一律に読むと強過ぎる。健康影響には物質の毒性・濃度・曝露期間が必要で、濃縮や不要という性質だけで被害は決まらない。 よって○・○・○・○の候補5を選ぶ。

\Needspace{4\baselineskip}\paragraph{選択肢の検討（査読待ち初稿）}\mbox{}\par

\noindent\textbf{1：整合しない} ×・○・○・○。個別判断と異なるのはアである。各記述の理由は確認条件欄に示す。\par

\noindent\textbf{2：整合しない} ○・×・○・○。個別判断と異なるのはイである。各記述の理由は確認条件欄に示す。\par

\noindent\textbf{3：整合しない} ○・○・×・○。個別判断と異なるのはウである。各記述の理由は確認条件欄に示す。\par

\noindent\textbf{4：整合しない} ○・○・○・×。個別判断と異なるのはエである。各記述の理由は確認条件欄に示す。\par

\noindent\textbf{5：整合する} ○・○・○・○。個別判断の○・○・○・○と全項目で一致する。\par

\paragraph{再利用する考え方} 活動・物質・制度の対象と評価の境界を固定し、用語の入替、全部/一部、法的義務/目標、出題時/現在を区別する。資源と排出の境界、濃縮と有害性の区別を個別の技術活動に結び付ける。

\paragraph{共通理論} \hyperref[theory:environment:carbon]{第\ref*{theory:environment:carbon}節}, \hyperref[theory:environment:development]{第\ref*{theory:environment:development}節}, \hyperref[theory:environment:zero]{第\ref*{theory:environment:zero}節}, \hyperref[theory:environment:accumulation]{第\ref*{theory:environment:accumulation}節}

\par 参考：\cite{env_offset2008,env_development,env_zero2008,env_accumulation}。

\paragraph{関連問題} \hyperref[q:r01-ret-apt-11]{R01-RET-APT-11}

\par 難易度 2（校準初稿）、優先度未評価。

\subsection{令和元年度 問6}
\label{q:r01-reg-apt-06}

\noindent\textbf{原典：} sources/R01/aptitude.pdf、PDF物理7頁。

\textbf{公式正答：} 1（single）。

\par\textit{原典本文の校正と独自解説は未作成。分類のみの登録項目。}

\subsection{令和元年度再試験 問4}
\label{q:r01-ret-apt-04}

\noindent\textbf{原典：} sources/R01\_retest/aptitude.pdf、PDF物理4頁。

\textbf{公式正答：} 1（single）。

\par\textit{原典本文の校正と独自解説は未作成。分類のみの登録項目。}

\subsection{令和3年度 問9}
\label{q:r03-reg-apt-09}

\noindent\textbf{原典：} sources/R03/aptitude.pdf、PDF物理9頁。

\textbf{公式正答：} 1（single）。

\par\textit{原典本文の校正と独自解説は未作成。分類のみの登録項目。}

\subsection{令和5年度 問10}
\label{q:r05-reg-apt-10}

\noindent\textbf{原典：} sources/R05/aptitude.pdf、PDF物理10頁。

\textbf{公式正答：} 3（single）。

\par\textit{原典本文の校正と独自解説は未作成。分類のみの登録項目。}

\subsection{令和5年度 問15}
\label{q:r05-reg-apt-15}

\noindent\textbf{原典：} sources/R05/aptitude.pdf、PDF物理15頁。

\textbf{公式正答：} 3（single）。

\paragraph{問題の要旨} 環境基本法の典型7公害として列挙されないものを選ぶ。

\Needspace{4\baselineskip}\paragraph{確認した条件・記述（独自要約）}\mbox{}\par

\noindent\textbf{方向：} 列挙されないもの\par

\paragraph{問われる核心} 環境法の公害類型と行為・被害の区別

\paragraph{必要知識} 環境法の公害類型と行為・被害の区別。7項目は大気汚染・水質汚濁・土壌汚染・騒音・振動・地盤沈下・悪臭である。3は原因となる行為の名前であり列挙された作用の種類ではない。法の限定列挙と日常語の広い環境問題を区別する。

\paragraph{考え方・正答を選ぶ理由（査読待ち初稿）} 7項目は大気汚染・水質汚濁・土壌汚染・騒音・振動・地盤沈下・悪臭である。3は原因となる行為の名前であり列挙された作用の種類ではない。法の限定列挙と日常語の広い環境問題を区別する。

\Needspace{4\baselineskip}\paragraph{選択肢の検討（査読待ち初稿）}\mbox{}\par

\noindent\textbf{1：列挙される} 騒音。第2条3項に列挙される。\par

\noindent\textbf{2：条件付きで列挙される} 地盤の沈下。法は鉱物掘採のための土地掘削によるものを除く条件で列挙する。\par

\noindent\textbf{3：列挙されない} 廃棄物投棄。それ自体は7項目の列挙名称ではない。投棄が水質・土壌等の公害を生む場合や別法違反となることを否定しない。\par

\noindent\textbf{4：列挙される} 悪臭。第2条3項に列挙される。\par

\noindent\textbf{5：列挙される} 振動。第2条3項に列挙される。\par

\paragraph{再利用する考え方} 活動・物質・制度の対象と評価の境界を固定し、用語の入替、全部/一部、法的義務/目標、出題時/現在を区別する。環境法の公害類型と行為・被害の区別を個別の技術活動に結び付ける。

\paragraph{共通理論} \hyperref[theory:environment:pollution]{第\ref*{theory:environment:pollution}節}

\par 参考：\cite{env_pollution}。

\par 難易度 2（校準初稿）、優先度未評価。

\subsection{令和5年度 問5}
\label{q:r05-reg-apt-05}

\noindent\textbf{原典：} sources/R05/aptitude.pdf、PDF物理5頁。

\textbf{公式正答：} 1（single）。

\par\textit{原典本文の校正と独自解説は未作成。分類のみの登録項目。}

\subsection{令和7年度 問4}
\label{q:r07-reg-apt-04}

\noindent\textbf{原典：} sources/R07/aptitude.pdf、PDF物理4頁。

\textbf{公式正答：} 3（single）。

\par\textit{原典本文の校正と独自解説は未作成。分類のみの登録項目。}

\subsection{令和7年度 問6}
\label{q:r07-reg-apt-06}

\noindent\textbf{原典：} sources/R07/aptitude.pdf、PDF物理6頁。

\textbf{公式正答：} 5（single）。

\par\textit{原典本文の校正と独自解説は未作成。分類のみの登録項目。}

\section{安全・リスク評価}

\subsection{令和5年度 問13}
\label{q:r05-reg-apt-13}

\noindent\textbf{原典：} sources/R05/aptitude.pdf、PDF物理13頁。

\textbf{公式正答：} 5（single）。

\par\textit{原典本文の校正と独自解説は未作成。分類のみの登録項目。}

\section{再生可能資源・熱移送と二次エネルギー}

\subsection{令和3年度 問11}
\label{q:r03-reg-apt-11}

\noindent\textbf{原典：} sources/R03/aptitude.pdf、PDF物理11頁。

\textbf{公式正答：} 3（single）。

\par\textbf{検討未確定：} 公式3・初回3を維持。問題幹の再生可能エネルギーは温室効果ガスを排出しないとの全称は、バイオガス燃焼やライフサイクルを含めれば不正確。パリ協定の採択2015/発効2016を区別。個別5記述の資源分類と幹の一般化を分離。

\paragraph{問題の要旨} 石炭、空気熱、水素、潮汐、バイオガスを資源の再生性と媒体で分類する。

\Needspace{4\baselineskip}\paragraph{確認した条件・記述（独自要約）}\mbox{}\par

\noindent\textbf{方向：} 個々の記述の○/×に適合する組合せ。\par

\noindent\textbf{並び：} ア・イ・ウ・エ・オ\par

\noindent\textbf{ア：} 石炭も植物から再び作られるため再生可能。 判定×：地質時間の形成は人間の利用時間で補充される再生可能性に当たらない。\par

\noindent\textbf{イ：} 空気熱をヒートポンプで温熱・冷熱利用。 判定○：自然界の熱を仕事を加えて移す技術である。温熱ではQ高=Q低+Wで、自然熱利用は入力仕事ゼロを意味しない。\par

\noindent\textbf{ウ：} 水素燃料は各種原料から製造するのでそれ自体は再生可能資源でない。 判定○：水素は二次エネルギー媒体で、製造元により再生可能性・排出が異なる。天然水素が絶対存在しないとはしない。\par

\noindent\textbf{エ：} 周期的潮汐のエネルギーを発電に利用。 判定○：人間の利用時間で周期的に補充される自然の流れを使う再生可能な供給である。\par

\noindent\textbf{オ：} 微生物分解で得るバイオガスは再生可能でない。 判定×：持続的に補充される生物資源由来のメタンはバイオマスエネルギーに含まれる。\par

\paragraph{問われる核心} 資源の再生時間、熱移送とエネルギー媒体

\paragraph{必要知識} 資源の再生時間、熱移送とエネルギー媒体。アは×。地質時間の形成は人間の利用時間で補充される再生可能性に当たらない。 イは○。自然界の熱を仕事を加えて移す技術である。温熱ではQ高=Q低+Wで、自然熱利用は入力仕事ゼロを意味しない。 ウは○。水素は二次エネルギー媒体で、製造元により再生可能性・排出が異なる。天然水素が絶対存在しないとはしない。 エは○。人間の利用時間で周期的に補充される自然の流れを使う再生可能な供給である。 オは×。持続的に補充される生物資源由来のメタンはバイオマスエネルギーに含まれる。 よって×・○・○・○・×の候補3を選ぶ。

\paragraph{考え方・正答を選ぶ理由（査読待ち初稿）} アは×。地質時間の形成は人間の利用時間で補充される再生可能性に当たらない。 イは○。自然界の熱を仕事を加えて移す技術である。温熱ではQ高=Q低+Wで、自然熱利用は入力仕事ゼロを意味しない。 ウは○。水素は二次エネルギー媒体で、製造元により再生可能性・排出が異なる。天然水素が絶対存在しないとはしない。 エは○。人間の利用時間で周期的に補充される自然の流れを使う再生可能な供給である。 オは×。持続的に補充される生物資源由来のメタンはバイオマスエネルギーに含まれる。 よって×・○・○・○・×の候補3を選ぶ。

\Needspace{4\baselineskip}\paragraph{選択肢の検討（査読待ち初稿）}\mbox{}\par

\noindent\textbf{1：整合しない} ○・○・○・○・○。個別判断と異なるのはア・オである。各記述の理由は確認条件欄に示す。\par

\noindent\textbf{2：整合しない} ○・×・○・×・○。個別判断と異なるのはア・イ・エ・オである。各記述の理由は確認条件欄に示す。\par

\noindent\textbf{3：整合する} ×・○・○・○・×。個別判断の×・○・○・○・×と全項目で一致する。\par

\noindent\textbf{4：整合しない} ×・○・×・○・×。個別判断と異なるのはウである。各記述の理由は確認条件欄に示す。\par

\noindent\textbf{5：整合しない} ×・×・×・×・○。個別判断と異なるのはイ・ウ・エ・オである。各記述の理由は確認条件欄に示す。\par

\paragraph{再利用する考え方} 活動・物質・制度の対象と評価の境界を固定し、用語の入替、全部/一部、法的義務/目標、出題時/現在を区別する。資源の再生時間、熱移送とエネルギー媒体を個別の技術活動に結び付ける。

\paragraph{共通理論} \hyperref[theory:environment:energy]{第\ref*{theory:environment:energy}節}

\par 参考：\cite{env_renewable2022,env_hydrogen2015,env_biomass,env_paris}。

\par 難易度 2（校準初稿）、優先度未評価。

\section{SDGsの普遍性・統合・人間の安全保障}

\subsection{令和元年度 問15}
\label{q:r01-reg-apt-15}

\noindent\textbf{原典：} sources/R01/aptitude.pdf、PDF物理14頁。

\textbf{公式正答：} 5（single）。

\paragraph{問題の要旨} SDGsの対象、原則、年限、目標数、参加者、範囲、評価の7記述を数える。

\Needspace{4\baselineskip}\paragraph{確認した条件・記述（独自要約）}\mbox{}\par

\noindent\textbf{方向：} 正しい記述の数。\par

\noindent\textbf{並び：} ア・イ・ウ・エ・オ・カ・キ\par

\noindent\textbf{ア：} 開発途上国だけの目標。 ×：先進国を含むすべての国へ適用する普遍性を持つ。\par

\noindent\textbf{イ：} 普遍性・包摂性・参画型・統合性・透明性。 ○：日本の実施指針が掲げる実施原則に一致し、最後は透明性と説明責任である。\par

\noindent\textbf{ウ：} 2030年を年限とする。 ○：2030アジェンダの年限である。個別ターゲットには2020年など異なる期限もある。\par

\noindent\textbf{エ：} 17の国際目標。 ○：目標は17、ターゲット169、指標はそれらと別の測定体系。\par

\noindent\textbf{オ：} 日本の取組は大企業・業界団体だけ。 ×：自治体・市民・個人等も参画し企業規模で排除しない。\par

\noindent\textbf{カ：} 環境問題に特化。 ×：経済・社会・環境を統合し、貧困・教育・平等なども含む。\par

\noindent\textbf{キ：} 指標と定期フォローアップ・評価公表。 ○：目標の達成を測る説明責任とレビューを必要とする。\par

\paragraph{問われる核心} SDGsの普遍性・統合性と説明責任

\paragraph{必要知識} SDGsの普遍性・統合性と説明責任。ア×・イ○・ウ○・エ○・オ×・カ×・キ○で、正しい4件、候補5。個々の理由は確認条件欄へ記録。2030期限と個別ターゲットの期限、目標/ターゲット/指標の階層を混同しない。

\paragraph{考え方・正答を選ぶ理由（査読待ち初稿）} ア×・イ○・ウ○・エ○・オ×・カ×・キ○で、正しい4件、候補5。個々の理由は確認条件欄へ記録。2030期限と個別ターゲットの期限、目標/ターゲット/指標の階層を混同しない。

\Needspace{4\baselineskip}\paragraph{選択肢の検討（査読待ち初稿）}\mbox{}\par

\noindent\textbf{1：整合しない} 0。イウエキの4件を全て除外している。\par

\noindent\textbf{2：整合しない} 1。4件中3件を取り落としている。\par

\noindent\textbf{3：整合しない} 2。4件中2件を取り落としている。\par

\noindent\textbf{4：整合しない} 3。4件中1件を取り落としている。\par

\noindent\textbf{5：整合する} 4。正しいイウエキの4件と一致。\par

\paragraph{再利用する考え方} 活動・物質・制度の対象と評価の境界を固定し、用語の入替、全部/一部、法的義務/目標、出題時/現在を区別する。SDGsの普遍性・統合性と説明責任を個別の技術活動に結び付ける。

\paragraph{共通理論} \hyperref[theory:environment:sdgs]{第\ref*{theory:environment:sdgs}節}, \hyperref[theory:environment:implementation]{第\ref*{theory:environment:implementation}節}

\par 参考：\cite{env_sdgs_agenda,env_sdgs2016,env_sdgs_un}。

\paragraph{関連問題} \hyperref[q:r04-reg-apt-14]{R04-REG-APT-14}

\par 難易度 2（校準初稿）、優先度未評価。

\subsection{令和元年度再試験 問14}
\label{q:r01-ret-apt-14}

\noindent\textbf{原典：} sources/R01\_retest/aptitude.pdf、PDF物理15頁。

\textbf{公式正答：} ALL（all\_awarded）。

\par\textbf{独立解答と公式正答に相違あり：} 独立解答 3。原典候補3の感梁症という誤植で的確な選択肢なし、公式全員得点。初回3は語を見落とした判断として不変保存。学習内容の指針照合と採点措置を分離。

\paragraph{問題の要旨} 2016年SDGs実施指針の4空欄と、原典選択肢の誤植による全員得点を記録する。

\Needspace{4\baselineskip}\paragraph{確認した条件・記述（独自要約）}\mbox{}\par

\noindent\textbf{空欄ア：} 地球規模の連鎖課題：感染症。\par

\noindent\textbf{空欄イ：} 従来の開発課題：開発途上国。\par

\noindent\textbf{空欄ウ：} 経済・社会と並ぶ環境。\par

\noindent\textbf{空欄エ：} 17目標と169ターゲット。\par

\noindent\textbf{原典：} 候補3の感梁症を感染症へ無断で校正しない。\par

\noindent\textbf{採点：} 公式正答ALL、全員得点。初回独立解答3。\par

\paragraph{問われる核心} SDGs実施指針の統合・普遍性と誤植の採点措置

\paragraph{必要知識} SDGs実施指針の統合・普遍性と誤植の採点措置。指針原文は感染症・開発途上国・環境・17となる。初回独立解答は3を保存したが、原本候補3には感梁症と印刷されている。公式注記は正答としていた候補の一部に誤植があり的確な選択肢がなく、受験者全員に得点を与える措置である。公式ALLを維持し、学習用に正しい語を示すことと原典の改変を区別する。

\paragraph{考え方・正答を選ぶ理由（査読待ち初稿）} 指針原文は感染症・開発途上国・環境・17となる。初回独立解答は3を保存したが、原本候補3には感梁症と印刷されている。公式注記は正答としていた候補の一部に誤植があり的確な選択肢がなく、受験者全員に得点を与える措置である。公式ALLを維持し、学習用に正しい語を示すことと原典の改変を区別する。

\Needspace{4\baselineskip}\paragraph{選択肢の検討（査読待ち初稿）}\mbox{}\par

\noindent\textbf{1：意図と不整合・全員得点} 国際紛争・先進国・環境・15。指針の感染症・開発途上国・環境・17とア・イ・エが異なる。\par

\noindent\textbf{2：意図と不整合・全員得点} 感染症・先進国・教育・15。イ・ウ・エが異なる。\par

\noindent\textbf{3：原典誤植・全員得点} 感梁症・開発途上国・環境・17（原典の字）。意図上近いが原典は感梁症という誤植であり、的確な候補でない。初回は感染症として3を選んだが、この見落としを後続の原本再照合で発見した。\par

\noindent\textbf{4：意図と不整合・全員得点} 国際紛争・開発途上国・教育・17。ア・ウが異なる。\par

\noindent\textbf{5：意図と不整合・全員得点} 感染症・開発途上国・教育・17。ウの教育が三側面の環境と異なる。\par

\paragraph{再利用する考え方} 活動・物質・制度の対象と評価の境界を固定し、用語の入替、全部/一部、法的義務/目標、出題時/現在を区別する。SDGs実施指針の統合・普遍性と誤植の採点措置を個別の技術活動に結び付ける。

\paragraph{共通理論} \hyperref[theory:environment:sdgs]{第\ref*{theory:environment:sdgs}節}, \hyperref[theory:environment:implementation]{第\ref*{theory:environment:implementation}節}

\par 参考：\cite{env_sdgs2016}。

\paragraph{関連問題} \hyperref[q:r06-reg-apt-06]{R06-REG-APT-06}

\par 難易度 2（校準初稿）、優先度未評価。

\subsection{令和3年度 問5}
\label{q:r03-reg-apt-05}

\noindent\textbf{原典：} sources/R03/aptitude.pdf、PDF物理5頁。

\textbf{公式正答：} 4（single）。

\paragraph{問題の要旨} SDGsの参加者、法的拘束力、背景、未来からの計画を判定する。

\Needspace{4\baselineskip}\paragraph{確認した条件・記述（独自要約）}\mbox{}\par

\noindent\textbf{方向：} 個々の記述の○/×に適合する組合せ。\par

\noindent\textbf{並び：} ア・イ・ウ・エ\par

\noindent\textbf{ア：} すべての参加者による誰一人取り残さない枠組み。 判定○：全ての国とstakeholderを対象とする。\par

\noindent\textbf{イ：} SDGsそのものに法的拘束力があり罰則対象。 判定×：SDGsの政治的目標を、それ自体罰則付きの条約義務としない。関連する国内法の遵守義務は別。\par

\noindent\textbf{ウ：} 環境・格差・紛争など将来への危機感を背景とする。 判定○：経済社会環境の相互作用を扱う課題設定である。\par

\noindent\textbf{エ：} バックキャストでなく現在からのフォーキャストが重要。 判定×：目標の将来像から現在へ必要条件を逆に求める方法が目標実現を設計する。実現可能性の分析を不要とするのでなく、将来像を現状延長だけで置き換えない。\par

\paragraph{問われる核心} SDGsの目標と法的義務、将来像からの計画

\paragraph{必要知識} SDGsの目標と法的義務、将来像からの計画。アは○。全ての国とstakeholderを対象とする。 イは×。SDGsの政治的目標を、それ自体罰則付きの条約義務としない。関連する国内法の遵守義務は別。 ウは○。経済社会環境の相互作用を扱う課題設定である。 エは×。目標の将来像から現在へ必要条件を逆に求める方法が目標実現を設計する。実現可能性の分析を不要とするのでなく、将来像を現状延長だけで置き換えない。 よって○・×・○・×の候補4を選ぶ。

\paragraph{考え方・正答を選ぶ理由（査読待ち初稿）} アは○。全ての国とstakeholderを対象とする。 イは×。SDGsの政治的目標を、それ自体罰則付きの条約義務としない。関連する国内法の遵守義務は別。 ウは○。経済社会環境の相互作用を扱う課題設定である。 エは×。目標の将来像から現在へ必要条件を逆に求める方法が目標実現を設計する。実現可能性の分析を不要とするのでなく、将来像を現状延長だけで置き換えない。 よって○・×・○・×の候補4を選ぶ。

\Needspace{4\baselineskip}\paragraph{選択肢の検討（査読待ち初稿）}\mbox{}\par

\noindent\textbf{1：整合しない} ○・×・○・○。個別判断と異なるのはエである。各記述の理由は確認条件欄に示す。\par

\noindent\textbf{2：整合しない} ○・○・○・×。個別判断と異なるのはイである。各記述の理由は確認条件欄に示す。\par

\noindent\textbf{3：整合しない} ×・○・×・○。個別判断と異なるのはア・イ・ウ・エである。各記述の理由は確認条件欄に示す。\par

\noindent\textbf{4：整合する} ○・×・○・×。個別判断の○・×・○・×と全項目で一致する。\par

\noindent\textbf{5：整合しない} ×・×・○・○。個別判断と異なるのはア・エである。各記述の理由は確認条件欄に示す。\par

\paragraph{再利用する考え方} 活動・物質・制度の対象と評価の境界を固定し、用語の入替、全部/一部、法的義務/目標、出題時/現在を区別する。SDGsの目標と法的義務、将来像からの計画を個別の技術活動に結び付ける。

\paragraph{共通理論} \hyperref[theory:environment:sdgs]{第\ref*{theory:environment:sdgs}節}, \hyperref[theory:environment:planning]{第\ref*{theory:environment:planning}節}

\par 参考：\cite{env_sdgs_agenda,env_sdgs_binding,env_sdgs2016}。

\par 難易度 2（校準初稿）、優先度未評価。

\subsection{令和4年度 問14}
\label{q:r04-reg-apt-14}

\noindent\textbf{原典：} sources/R04/aptitude.pdf、PDF物理13頁。

\textbf{公式正答：} 3（single）。

\paragraph{問題の要旨} SDGsの対象国、MDGsとの継承、3側面、人権と参加の7記述を判定する。

\Needspace{4\baselineskip}\paragraph{確認した条件・記述（独自要約）}\mbox{}\par

\noindent\textbf{方向：} 個々の記述の○/×に適合する組合せ。\par

\noindent\textbf{並び：} ア・イ・ウ・エ・オ・カ・キ\par

\noindent\textbf{ア：} 先進国だけが実行する。 判定×：先進国と途上国の双方を含む普遍性に反する。\par

\noindent\textbf{イ：} MDGsの未達成課題を引き継いで全うする。 判定○：前文のMDGsを基に残課題を完成する趣旨。\par

\noindent\textbf{ウ：} 経済・社会・環境を調和させる。 判定○：持続可能な開発の3側面を統合する。\par

\noindent\textbf{エ：} 誰一人取り残さない。 判定○：脆弱な人々と不平等へ焦点を当てる包摂性。\par

\noindent\textbf{オ：} 人権とジェンダー平等・女性女児の能力強化。 判定○：前文と目標5、段落20に対応。\par

\noindent\textbf{カ：} 全stakeholderが協同的パートナーシップで実行。 判定○：政府だけに限定しない。\par

\noindent\textbf{キ：} 環境問題に特化する。 判定×：教育・貧困・平等・経済等も含むため限定が誤り。\par

\paragraph{問われる核心} SDGsの普遍・包摂・統合と協働

\paragraph{必要知識} SDGsの普遍・包摂・統合と協働。アは×。先進国と途上国の双方を含む普遍性に反する。 イは○。前文のMDGsを基に残課題を完成する趣旨。 ウは○。持続可能な開発の3側面を統合する。 エは○。脆弱な人々と不平等へ焦点を当てる包摂性。 オは○。前文と目標5、段落20に対応。 カは○。政府だけに限定しない。 キは×。教育・貧困・平等・経済等も含むため限定が誤り。 よって×・○・○・○・○・○・×の候補3を選ぶ。

\paragraph{考え方・正答を選ぶ理由（査読待ち初稿）} アは×。先進国と途上国の双方を含む普遍性に反する。 イは○。前文のMDGsを基に残課題を完成する趣旨。 ウは○。持続可能な開発の3側面を統合する。 エは○。脆弱な人々と不平等へ焦点を当てる包摂性。 オは○。前文と目標5、段落20に対応。 カは○。政府だけに限定しない。 キは×。教育・貧困・平等・経済等も含むため限定が誤り。 よって×・○・○・○・○・○・×の候補3を選ぶ。

\Needspace{4\baselineskip}\paragraph{選択肢の検討（査読待ち初稿）}\mbox{}\par

\noindent\textbf{1：整合しない} ×・×・○・○・○・○・○。個別判断と異なるのはイ・キである。各記述の理由は確認条件欄に示す。\par

\noindent\textbf{2：整合しない} ×・○・×・○・×・○・×。個別判断と異なるのはウ・オである。各記述の理由は確認条件欄に示す。\par

\noindent\textbf{3：整合する} ×・○・○・○・○・○・×。個別判断の×・○・○・○・○・○・×と全項目で一致する。\par

\noindent\textbf{4：整合しない} ○・×・○・×・○・×・○。個別判断と異なるのはア・イ・エ・カ・キである。各記述の理由は確認条件欄に示す。\par

\noindent\textbf{5：整合しない} ×・○・○・○・○・×・×。個別判断と異なるのはカである。各記述の理由は確認条件欄に示す。\par

\paragraph{再利用する考え方} 活動・物質・制度の対象と評価の境界を固定し、用語の入替、全部/一部、法的義務/目標、出題時/現在を区別する。SDGsの普遍・包摂・統合と協働を個別の技術活動に結び付ける。

\paragraph{共通理論} \hyperref[theory:environment:sdgs]{第\ref*{theory:environment:sdgs}節}

\par 参考：\cite{env_sdgs_agenda,env_sdgs_un}。

\paragraph{関連問題} \hyperref[q:r01-reg-apt-15]{R01-REG-APT-15}

\par 難易度 2（校準初稿）、優先度未評価。

\subsection{令和6年度 問2}
\label{q:r06-reg-apt-02}

\noindent\textbf{原典：} sources/R06/aptitude.pdf、PDF物理2頁。

\textbf{公式正答：} 5（single）。

\paragraph{問題の要旨} 2030アジェンダのジェンダー段落の4空欄に入らない語を選ぶ。

\Needspace{4\baselineskip}\paragraph{確認した条件・記述（独自要約）}\mbox{}\par

\noindent\textbf{並び：} ア能力強化（エンパワーメント）・イ雇用・ウリーダーシップ・エ意思決定。\par

\noindent\textbf{方向：} 4空欄のいずれにも入らない語。\par

\paragraph{問われる核心} SDGsのジェンダー平等と指定文脈の能力・参画

\paragraph{必要知識} SDGsのジェンダー平等と指定文脈の能力・参画。段落20の並びは能力強化、雇用、リーダーシップ、意思決定である。よってケア5がどの空欄にも入らない。文脈の範囲を指定段落に固定し、他のターゲットに語があることと区別する。

\paragraph{考え方・正答を選ぶ理由（査読待ち初稿）} 段落20の並びは能力強化、雇用、リーダーシップ、意思決定である。よってケア5がどの空欄にも入らない。文脈の範囲を指定段落に固定し、他のターゲットに語があることと区別する。

\Needspace{4\baselineskip}\paragraph{選択肢の検討（査読待ち初稿）}\mbox{}\par

\noindent\textbf{1：イに対応} 雇用。イに入り、働く機会の平等を示す。\par

\noindent\textbf{2：ウに対応} リーダーシップ。ウに入り、指導的役割へ参画する機会。\par

\noindent\textbf{3：エに対応} 意思決定。エに入り、決定過程で同等の機会。\par

\noindent\textbf{4：アに対応} エンパワーメント。アに入り、女性女児の能力強化。外務省仮訳は能力強化と表す。\par

\noindent\textbf{5：どれにも入らない} ケア。この段落の4空欄には入らない。ただし目標5.4は無償のケア/家事労働を扱い、SDGs全体でケアが無関係という結論ではない。\par

\paragraph{再利用する考え方} 活動・物質・制度の対象と評価の境界を固定し、用語の入替、全部/一部、法的義務/目標、出題時/現在を区別する。SDGsのジェンダー平等と指定文脈の能力・参画を個別の技術活動に結び付ける。

\paragraph{共通理論} \hyperref[theory:environment:gender]{第\ref*{theory:environment:gender}節}

\par 参考：\cite{env_sdgs_agenda}。

\paragraph{関連問題} \hyperref[q:r06-reg-apt-06]{R06-REG-APT-06}

\par 難易度 2（校準初稿）、優先度未評価。

\subsection{令和6年度 問6}
\label{q:r06-reg-apt-06}

\noindent\textbf{原典：} sources/R06/aptitude.pdf、PDF物理6頁。

\textbf{公式正答：} 2（single）。

\paragraph{問題の要旨} SDGsの3側面、統合性、普遍性、人間の安全保障の4空欄に入らない語を選ぶ。

\Needspace{4\baselineskip}\paragraph{確認した条件・記述（独自要約）}\mbox{}\par

\noindent\textbf{方向：} 指定4空欄のどれにも入らない。\par

\noindent\textbf{並び：} ア経済・イ統合・ウ普遍・エ安全。\par

\paragraph{問われる核心} SDGsの統合・普遍性と人間の安全保障

\paragraph{必要知識} SDGsの統合・普遍性と人間の安全保障。ア経済、イ統合、ウ普遍、エ安全と対応し、候補2の平和が空欄にない。安全保障は国家間の軍事だけでなく、脆弱な一人ひとりが尊厳ある生活を送れるという人間中心の概念である。

\paragraph{考え方・正答を選ぶ理由（査読待ち初稿）} ア経済、イ統合、ウ普遍、エ安全と対応し、候補2の平和が空欄にない。安全保障は国家間の軍事だけでなく、脆弱な一人ひとりが尊厳ある生活を送れるという人間中心の概念である。

\Needspace{4\baselineskip}\paragraph{選択肢の検討（査読待ち初稿）}\mbox{}\par

\noindent\textbf{1：エに対応} 安全。エの人間の安全保障を構成する。\par

\noindent\textbf{2：空欄にない} 平和。この4空欄には入らない。平和はSDG16や前文の重要事項であり、SDGsの対象外という意味ではない。\par

\noindent\textbf{3：アに対応} 経済。アに入り社会・経済・環境の3側面を成す。\par

\noindent\textbf{4：ウに対応} 普遍。ウに入りすべての国の目標という普遍性。\par

\noindent\textbf{5：イに対応} 統合。イに入り相互関係のある目標を統合的に解決する。\par

\paragraph{再利用する考え方} 活動・物質・制度の対象と評価の境界を固定し、用語の入替、全部/一部、法的義務/目標、出題時/現在を区別する。SDGsの統合・普遍性と人間の安全保障を個別の技術活動に結び付ける。

\paragraph{共通理論} \hyperref[theory:environment:sdgs]{第\ref*{theory:environment:sdgs}節}, \hyperref[theory:environment:implementation]{第\ref*{theory:environment:implementation}節}

\par 参考：\cite{env_sdgs_pamphlet,env_humansecurity}。

\paragraph{関連問題} \hyperref[q:r01-ret-apt-14]{R01-RET-APT-14}, \hyperref[q:r06-reg-apt-02]{R06-REG-APT-02}

\par 難易度 2（校準初稿）、優先度未評価。

\section{社会的責任・安全の標準}

\subsection{令和6年度 問4}
\label{q:r06-reg-apt-04}

\noindent\textbf{原典：} sources/R06/aptitude.pdf、PDF物理4頁。

\textbf{公式正答：} 1（single）。

\par\textit{原典本文の校正と独自解説は未作成。分類のみの登録項目。}
